% Folder 33, batch 05 — archivists' pages 81–100.
% Pass: Opus 5.5 (claude-opus-5-5), 2026-10-03 — first pass, unchecked against the pages by a human.
%
% Only the odd pages (81, 83, …, 99) carry his manuscript. The even pages are
% typed versos: a typescript of EGA IV (§ 19, pages headed IV-1008 to IV-1013,
% some with marginal corrections in a hand that is not clearly his) and pages
% of English-language mathematical typescripts unrelated to these notes. They
% are skipped.
\documentclass[11pt,a4paper]{article}
\input{../preamble/grothendieck}

\folder{33}
\batch{5}
\pages{81}{100}
\foldertitle{SGA 1965. MS [Manuscrit] brouillon : notes manuscrites (s.d.).}
\dating{1965-[vers 1971]}
\watermark{Édition de démonstration}

\begin{document}

\page{81}
\note{la page commence au milieu d'un argument venu des feuillets précédents (homomorphismes de Gysin pour un morphisme $f$ lissifiable au-dessus d'une base) ; les notations $h$, $g$, $d_{X/Y}$, $T_{X/Y}$ y sont définies.}

Donc on trouve (sans condition de finitude sur $f$) des \emph{homomorphismes de Gysin}
\[
\mathbb{R}_!h\bigl(\mathbb{L}h^{*}F^{\cdot}\otimes T_{X/Y}[2d_{X/Y}]\bigr)\longrightarrow\mathbb{R}_!g\bigl(\mathbb{L}g^{*}(F^{\cdot})\bigr)
\]
\note{dans l'indice de $T$, un $Z$ est biffé et remplacé par $Y$.}
en particulier
\[
\mathbb{R}_!^{\,i+2d_{X/Y}}h\bigl(\mathbb{L}h^{*}(F^{\cdot})\bigr)\otimes T_{X/Y}\longrightarrow\mathbb{R}_!^{\,i}g\bigl(\mathbb{L}g^{*}(F^{\cdot})\bigr)
\]
\note{après $T_{X/Y}$, un crochet $[2d_{X/Y}]$ est biffé.}
et si $f$ est propre, globalement
\[
\mathbb{H}^{i+2d_{X/Y}}\bigl(X,\mathbb{L}h^{*}(F^{\cdot})\bigr)\otimes T_{X/Y}\longrightarrow\mathbb{H}^{i}\bigl(Y,\mathbb{L}g^{*}(F^{\cdot})\bigr)
\]
\note{même biffure du décalage $[2d_{X/Y}]$ dans cette formule.}

Nous admettrons que \ill{} \uncertain{satisfont} (\uncertain{figurant} dans le contexte 1°, et « \ill{} » dans le contexte 3°) à la condition de transitivité \uncertain{qui s'impose} \ill{}.

\page{83}
\textbf{24} \struck{\ill{}} \emph{Transitivités pour $\mathbb{R}^!f$, et homomorphismes de Gysin}

1°) \uncertain{L'homomorphisme} de \uncertain{Gysin}, \ill{} \uncertain{trace}, pour un morphisme $f\colon X\to Y$ localement compactifiable, \struck{et} et un $G^{\cdot}\in\operatorname{Ob}D^{+}(Y)$, on a l'homomorphisme naturel
\[
(*)\qquad \mathbb{R}_!f\,\mathbb{R}^!f(G^{\cdot})\longrightarrow G^{\cdot}
\]
d'où \ill{} homomorphismes qui s'en déduisent,
\[
\begin{cases}
\mathbb{R}^{i}_!f\bigl(\mathbb{R}^!f(G^{\cdot})\bigr)\longrightarrow\mathcal{H}^{i}(G^{\cdot})\\
\mathbb{H}^{i}_{\Phi(f)}\bigl(X,\mathbb{R}^!f(G^{\cdot})\bigr)\longrightarrow\mathbb{H}^{i}(Y,G^{\cdot})
\end{cases}
\]
\note{devant la seconde ligne, un premier début ($\mathbb{R}^i_!$ ?) est biffé ; l'indice de supports se lit $\Phi(f)$ sans certitude.}
et aussi, pour $f$ \uncertain{propre}, on \uncertain{a aussi} un
\[
\mathbb{H}^{i}\bigl(X,\mathbb{R}^!f(G^{\cdot})\bigr)\longrightarrow\mathbb{H}^{i}(Y,G^{\cdot})
\]

\emph{Variante} : si $X$, $Y$ sont sur $Z$, \uncertain{donc} on \uncertain{trouve}, en \uncertain{utilisant} $\mathbb{R}_!g$ : \ill{},

\begin{tikzcd}
X \arrow[r, "f"] \arrow[dr, "h"'] & Y \arrow[d, "g"] \\
& Z
\end{tikzcd}

\[
\text{d'où}\quad
\begin{cases}
\mathbb{R}_!h\bigl(\mathbb{R}^!f(G^{\cdot})\bigr)\longrightarrow\mathbb{R}_!g(G^{\cdot})\\
\mathbb{H}^{i}_{\Phi(h)}\bigl(X,\mathbb{R}^!f(G^{\cdot})\bigr)\longrightarrow\mathbb{H}^{i}_{\Phi(g)}(Y,G^{\cdot})
\end{cases}
\]
\note{devant la seconde ligne, un $\mathbb{R}^i$ est biffé.}

\page{85}
\textbf{25} \struck{\ill{}} \struck{\ill{} \ill{} \ill{} homomorphismes de} \emph{Formulaire de $\mathbb{R}_!f$, $\mathbb{R}f_{*}$, … $\mathbb{R}^!f$, $\mathbb{L}f^{*}$}
\note{le numéro est récrit, un 6 corrigé en 5 ; au-dessus du titre biffé, un mot lui-même biffé ; sous le titre, un mot biffé (\ill{}).}
et $\overset{\mathbb{L}}{\otimes}$ et $\mathbb{R}\mathcal{H}om$
\marginal{sous l'accolade qui réunit $\mathbb{R}_!f,\dots,\mathbb{L}f^{*}$ : opérations entre \uncertain{catégories} (opérations internes).}

$\mathbb{R}_!f$, $\mathbb{R}^!f$ définis seulement si $f$ compactifiable, resp.\ lissifiable (strictement, i.e.\ pas de problèmes relatifs \ill{} \uncertain{de dim. bornée}).

1°) \emph{Transitivités}
\[
\text{Trans.}*\ \begin{cases}
\mathbb{R}(gf)_{*}\simeq\mathbb{R}g_{*}\,\mathbb{R}f_{*}\\
\mathbb{L}(gf)^{*}\simeq\mathbb{L}f^{*}\,\mathbb{L}g^{*}
\end{cases}
\qquad
\text{Trans.}!\ \begin{cases}
\mathbb{R}_!(gf)\simeq\mathbb{R}_!g\,(\mathbb{R}_!f)\\
\mathbb{R}^!(gf)\simeq(\mathbb{R}^!f)(\mathbb{R}^!g)
\end{cases}
\]
\note{les étiquettes à gauche sont récrites sur des mots biffés.}

Formules de type universel. Formules de droite déjà vues pour $\mathbb{R}_!$. Pour $\mathbb{R}^!$, cela résulte « immédiatement » par adjonction. Une démonstration \uncertain{indépendante} \ill{} lorsque \uncertain{l'on suppose} que $Y$ se plonge dans $Y'$ lisse sur $Z$, de façon que $X$ soit $Y'$-lissifiable [par \ill{} \ill{} \ill{}

\begin{tikzcd}
X \arrow[r, hook] \arrow[d] & X'_{Y} \arrow[r, hook] \arrow[dl] & X' \arrow[dl] \\
Y \arrow[r, hook] \arrow[d] & Y' \arrow[dl, "\text{lisse}"] & \\
Z & &
\end{tikzcd}

\ill{} \ill{} \ill{} \ill{} \uncertain{stable par toutes les opérations}, \ill{} \ill{} \ill{} de base \ill{}. \uncertain{$S$ admettant} un \uncertain{faisceau inversible} \ill{}, \struck{\ill{}} \ill{} \ill{} propriétés sur $S$ ; donc pour tout $X\to Y$, $X$ se plonge \ill{} dans un $\mathbb{P}^{r}_{Y}$ …]
\marginal{\uncertain{peut-on} \ill{} \uncertain{démontrer} dans la \uncertain{formulation}}

\page{87}
L'isomorphisme de transitivité se décompose \uncertain{donc} en trois isomorphismes de transitivité particuliers : a) deux morphismes lisses, b) \emph{Deux} morphismes dont l'un \uncertain{immersion}, c) \uncertain{commutativité} des \uncertain{images inverses} dans le cas (?). Les cas a) b) sont triviaux, le cas c) est essentiellement le th.\ de changement de base lisse.

\emph{NB.} Dans tous ces isomorphismes de transitivité, il faudrait vérifier que quand il y en a trois composés, on a la compatibilité habituelle. — Nous l'admettrons.

\struck{\ill{}}

2°) \emph{Formules d'adjonction}

\begin{tikzcd}
X \arrow[d, "f"] & F^{\cdot} \\
Y & G^{\cdot}
\end{tikzcd}

\[
\boxed{\operatorname{Hom}\bigl(\mathbb{R}_!f(F_{\cdot}),G^{\cdot}\bigr)\simeq\operatorname{Hom}\bigl(F_{\cdot},\mathbb{R}^!f(G^{\cdot})\bigr)}
\]
\[
[F_{\cdot}\in\operatorname{Ob}D^{b}(X),\ G^{\cdot}\in\operatorname{Ob}D^{+}(Y)]
\]
\[
\boxed{\operatorname{Hom}\bigl(\mathbb{L}f^{*}(G_{\cdot}),F^{\cdot}\bigr)\simeq\operatorname{Hom}\bigl(G_{\cdot},\mathbb{R}f_{*}(F^{\cdot})\bigr)}
\]
\[
[F^{\cdot}\in\operatorname{Ob}D^{+}(X),\ G^{\cdot}\in\operatorname{Ob}D^{b}(Y)]
\]

\page{89}
\emph{Rem.} La deuxième formule est triviale, ($f$ \uncertain{plat} \ill{} \uncertain{qui simplifie} …) \struck{résulte}
\begin{align*}
\operatorname{Hom}\bigl(G^{\cdot},\mathbb{R}f_{*}(F^{\cdot})\bigr)
&\simeq H^{0}\Bigl(\operatorname{Hom}^{\cdot}\bigl(G^{\cdot},f_{*}(C^{\cdot}(F^{\cdot}))\bigr)\Bigr)\\
&\simeq H^{0}\Bigl(\operatorname{Hom}^{\cdot}\bigl(f^{*}G^{\cdot},C^{\cdot}(F^{\cdot})\bigr)\Bigr)\\
&\simeq\operatorname{Hom}\bigl(\mathbb{L}f^{*}G^{\cdot},F\bigr)
\end{align*}

3°) \emph{Formules de projection}
\[
\boxed{\mathbb{R}\mathcal{H}om\bigl(\mathbb{R}_!f(F_{\cdot}),G^{\cdot}\bigr)\simeq\mathbb{R}f_{*}\Bigl(\mathbb{R}\mathcal{H}om\bigl(\mathbb{R}^!f(F_{\cdot}),G^{\cdot}\bigr)\Bigr)}
\]
\note{ainsi sur la page : à droite on attendrait $\mathbb{R}\mathcal{H}om(F_{\cdot},\mathbb{R}^!f(G^{\cdot}))$.}
\[
\boxed{\mathbb{R}\mathcal{H}om\bigl(G_{\cdot},\mathbb{R}f_{*}(F^{\cdot})\bigr)\simeq\mathbb{R}f_{*}\Bigl(\mathbb{R}\mathcal{H}om\bigl(\mathbb{L}f^{*}(G_{\cdot}),F^{\cdot}\bigr)\Bigr)}
\]
\[
\boxed{G_{\cdot}\overset{\mathbb{L}}{\otimes}\mathbb{R}f_{*}(F_{\cdot})\simeq\mathbb{R}f_{*}\bigl(G_{\cdot}\overset{\mathbb{L}}{\otimes}F_{\cdot}\bigr)}
\]
\note{dans cette troisième formule, des $!$ et $*$ sont surchargés et en partie biffés ; la lecture $\mathbb{R}f_{*}$ des deux côtés est incertaine.}
\marginal{\struck{i.e.\ pour $\mathbb{R}_{*}f(F_{\cdot})$} si $G_{\cdot}$ loc.\ libre (peut se généraliser : $G_{\cdot}$ loc.\ \uncertain{constant} \ill{} 2°, 3°)}

\emph{NB} Les \struck{\ill{}} isom.\ \ill{} \uncertain{formules} \ill{} \uncertain{des homomorphismes canoniques}
\begin{align*}
\operatorname{Tr}_{f}(G^{\cdot})&\colon\ \mathbb{R}_!f\,\mathbb{R}^!f(G^{\cdot})\longrightarrow G^{\cdot}\\
\tau_{f}(G^{\cdot})&\colon\ G_{\cdot}\longrightarrow\mathbb{R}f_{*}\bigl(\mathbb{L}f^{*}(G_{\cdot})\bigr)
\end{align*}

On a des énoncés de compatibilité entre les isom.\ en 2° ou 3°, i.e.\ \uncertain{entre} les homom.\ $\operatorname{Tr}_{f}$ et $\tau_{f}$,
\marginal{dans la marge gauche, sous un hachurage de crayon, quelques mots tournés (« pour \ill{} … ») ne se lisent pas.}

\page{91}
et les isom.\ de transitivité de 1° ….

4°) \emph{Formules de changement de base} (\uncertain{analogues} \uncertain{après} 1°)

\begin{tikzcd}
X \arrow[d, "f"] & X' \arrow[l, "g'"'] \arrow[d, "f'"] \\
Y & Y' \arrow[l, "g"]
\end{tikzcd}

\[
\boxed{\mathbb{L}g^{*}\,\mathbb{R}_!f\xrightarrow{\ \sim\ }\mathbb{R}_!f'\,\mathbb{L}g'^{*}}
\]
\marginal{changement de base pour morphisme propre}
\[
\begin{cases}
\mathbb{L}g^{*}\,\mathbb{R}f_{*}\xrightarrow{\ \sim\ }\mathbb{R}f'_{*}\,\mathbb{L}g'^{*} & \text{si $f$ propre [ou $g$ lisse]}\\
\mathbb{R}^!f'\,\mathbb{L}g^{*}\xleftarrow{\ \sim\ }\mathbb{L}g'^{*}\,\mathbb{R}^!f & \text{si $f$ ou $g$ lisse}
\end{cases}
\]
\note{accolade reliée à « $g$ \uncertain{quelconque} » dans la marge ; dans la seconde ligne, les accents et les $*$ sont surchargés, et la place des primes sur $f$ est incertaine.}

\[
\text{\struck{$\mathbb{R}^!g\,\mathbb{R}_!f\Rightarrow\mathbb{R}_!f'\,\mathbb{R}^!g'$}}
\]
\note{encadré biffé d'un trait ondulé ; à droite : « si $f$ \ill{} ».}
\[
\boxed{\mathbb{R}^!g\,\mathbb{R}f_{*}\simeq\mathbb{R}f'_{*}\,\mathbb{R}^!g'}
\]
\marginal{$g$ lissifiable … est une conséquence du précédent}
\note{une flèche va de « $g$ lissifiable » à l'encadré, une autre de l'encadré à « est une conséquence du précédent ».}

\marginal{si $f$ propre [ou $g$ lisse] ; si $f$ ou $g$ lisse $\to$ changement de base lisse ; (se décompose en cas lisse = changement de base lisse, et cas d'une immersion, trivial)}

5°) \emph{Formules d'induction}

\begin{tikzcd}
X \arrow[d, "f"] & \\
Y & F,G
\end{tikzcd}

\[
\boxed{\mathbb{R}^!f\bigl(\mathbb{R}\mathcal{H}om(F_{\cdot},G^{\cdot})\bigr)\simeq\mathbb{R}\mathcal{H}om\bigl(\mathbb{L}f^{*}(F_{\cdot}),\mathbb{R}^!f(G^{\cdot})\bigr)}
\]
\[
\boxed{\mathbb{L}f^{*}\bigl(F_{\cdot}\overset{\mathbb{L}}{\otimes}G_{\cdot}\bigr)\simeq\mathbb{L}f^{*}(F_{\cdot})\overset{\mathbb{L}}{\otimes}\mathbb{L}f^{*}(G_{\cdot})}
\]
\marginal{formules pour topos annelés quelconques …}

\page{93}
Établissons la première formule d'induction. On distingue les cas $f$ \emph{lisse}, et $f$ une immersion fermée.

a) \emph{Cas $f$ lisse}. La formule équivaut à :
\[
\mathbb{L}f^{*}\bigl(\mathbb{R}\mathcal{H}om(F_{\cdot},G^{\cdot})\bigr)\simeq\mathbb{R}\mathcal{H}om\bigl(\mathbb{L}f^{*}F_{\cdot},\mathbb{L}f^{*}G^{\cdot}\bigr)
\]
i.e.\ au fait que si $f$ est lisse, $\mathbb{L}f^{*}$ \emph{commute aux $\mathcal{E}xt^{i}$ locaux}, (si \ill{} \ill{}) \struck{et \ill{} cohomologique constructible}. On a \ill{} en tout cas un \ill{} \uncertain{flèche} \uncertain{évidente} $\to$, \ill{} à prouver que c'est un isomorphisme.

\struck{C'est local, \ill{} pour supposer $Y$, $X$ affines, $F_{\cdot}$ réduit \ill{} degré $0$, $G^{\cdot}$ \ill{} et injectif. Alors $F$ admet une filtration \ill{} des $F'$ loc.\ constant sur $Y'$ \ill{}, \ill{} $Y'\to Y$ immersion, \ill{} On peut supposer que $F=\mathbb{R}_!i(F')$ \ill{} On a une \uncertain{résolution} de $F$ \ill{} $Y'\to Y$ \ill{} $(\mathbb{R}_!i(A_{Y'}))$ \ill{}}
\note{tout ce passage, de « C'est local » à la fin de la page, est biffé de grandes croix et partiellement encadré ; on n'en donne que ce qui se lit. Au milieu, sous la biffure, « $G^{\cdot}$ \uncertain{tous} et injectif » est récrit.}

\page{95}
\begin{tikzcd}
X' \arrow[r, hook, "j"] \arrow[d, "f'"'] & X \arrow[d, "f"] \\
Y' \arrow[r, hook, "i"'] & Y
\end{tikzcd}

$F'$ \qquad $f^{*}(F)=j_!\bigl(f'^{*}(F')\bigr)$, où $f'^{*}(F')=F'_{X'}$.

\begin{align*}
\mathbb{R}\mathcal{H}om\bigl(\mathbb{R}i_!(F'),G^{\cdot}\bigr)&\simeq\mathbb{R}i_{*}\Bigl(\mathbb{R}\mathcal{H}om\bigl(F',\mathbb{R}i^!(G^{\cdot})\bigr)\Bigr)
&&\text{(dualité pour $i$)}\\
\mathbb{L}f^{*}(\ \cdots\ )&\simeq\mathbb{L}f^{*}\,\mathbb{R}i_{*}(\ \cdots\ )\\
&\simeq\mathbb{R}j_{*}\,\mathbb{L}f'^{*}\Bigl(\mathbb{R}\mathcal{H}om\bigl(F',\mathbb{R}i^!(G^{\cdot})\bigr)\Bigr)
&&\text{(chang.\ de base lisse)}\\
&\simeq\mathbb{R}j_{*}\Bigl(\mathbb{R}\mathcal{H}om\bigl(f'^{*}(F'),f'^{*}(\mathbb{R}i^!(G^{\cdot}))\bigr)\Bigr)\\
&\simeq\mathbb{R}j_{*}\Bigl(\mathbb{R}\mathcal{H}om\bigl(F'_{X'},\mathbb{R}j^!G^{\cdot}_{X}\bigr)\Bigr)\\
&\simeq\mathbb{R}\mathcal{H}om\bigl(j_!(F'_{X'}),G^{\cdot}_{X}\bigr)
&&\text{(dualité pour $j$)}\\
&\simeq\mathbb{R}\mathcal{H}om\bigl(f^{*}(F_{X}),f^{*}(G)\bigr)
\end{align*}
\note{commentaires de la page aux deux premières étapes de la chaîne : sous l'accolade de $\mathbb{R}\mathcal{H}om(F',\mathbb{R}i^!(G^{\cdot}))$, « \uncertain{supposition} \ill{} \uncertain{instant} : \uncertain{le th.} \uncertain{vrai} pour $f'\colon X'\to Y'$ et $F'$, $\mathbb{R}i^!G$ » ; sous $f'^{*}(\mathbb{R}i^!(G^{\cdot}))$, « $\simeq$ — th.\ de changement de base lisse : $\mathbb{R}j^!(f'^{*}(G^{\cdot}))$ » ; sous $j_!(F'_{X'})$, l'accolade porte $f^{*}(F_{X})$. Un mot à droite de la dernière ligne est biffé.}

On est \uncertain{ramené} donc au cas où $F$ est loc.\ constant. On peut donc (\uncertain{le th.} étant local)

\page{97}
supposer $F$ \emph{constant}, puis (\uncertain{le th.} \ill{} \uncertain{additif} en $A_{Y}^{n}$) $F=A_{Y}$, mais alors le th.\ est trivial.

b) \emph{Cas $f$ une immersion}. Alors le th.\ est essentiellement trivial, sans hypothèse de finitude, \ill{} \uncertain{ramené} au cas d'une immersion fermée. Alors
\begin{align*}
\mathbb{R}^!f\bigl(\mathbb{R}\mathcal{H}om(F_{\cdot},G^{\cdot})\bigr)&\simeq f^!\,\mathcal{H}om^{\cdot}\bigl(F_{\cdot},C^{\cdot}(G^{\cdot})\bigr)\\
&\simeq\mathcal{H}om^{\cdot}\bigl(f^{*}F_{\cdot},f^!C^{\cdot}(G^{\cdot})\bigr)\simeq\mathbb{R}\mathcal{H}om\bigl(\mathbb{L}f^{*}(F_{\cdot}),\mathbb{R}f^!(G^{\cdot})\bigr)
\end{align*}
\note{sous $\mathcal{H}om^{\cdot}(F_{\cdot},C^{\cdot}(G^{\cdot}))$ : « N.B.\ c'est flasque, \uncertain{pour} $f^!$ \uncertain{quelconque} » ; sous $f^!C^{\cdot}(G^{\cdot})$ : « N.B.\ c'est injectif ». Dans la première ligne, un $\mathbb{R}$ est surchargé ($\mathbb{R}^!$ récrit sur $\mathbb{R}$ ?), et après $f^!$ quelques lettres sont biffées.}

Les formules de projection, d'induction \uncertain{s'interprètent} comme des formules \struck{de \ill{}} \emph{d'échange}
\[
\begin{cases}
\mathbb{R}f_{*}\,D_{X}\simeq D_{Y}\,\mathbb{R}_!f\\
\mathbb{R}^!f\,D_{Y}\simeq D_{X}\,\mathbb{L}f^{*}
\end{cases}
\]
i.e.\ $D$ « échange » $\mathbb{R}f_{*}$ et $\mathbb{R}_!f$, $\mathbb{L}f^{*}$ et $\mathbb{R}^!f$.

\page{99}
\textbf{(26)} \emph{Formulaire faisant intervenir $\overset{\mathbb{L}}{\otimes}$ et $D^{-}$}

(1) Généralités valables sur tout topos annelé.

Module \emph{plat} $F$ : $F\otimes\cdot$ est exact

$F$ plat $\Longleftrightarrow$ Pour tout injectif $I$, $\mathcal{H}om(F,I)$ est injectif
\quad \emph{NB} Dans le cas actuel, il suffit de vérifier sur les fibres.

Tout Module $F$ \struck{admet des \ill{}} est quotient d'un plat, donc admet des résolutions plates : on prend les $A_{U}$ ($U\in\operatorname{Ob}X_{\text{ét}}$). Il s'ensuit \add{formellement} que pour tout $L_{\cdot}\in\operatorname{Ob}D^{-}(X)$, existe homom.\ $L'_{\cdot}\to L_{\cdot}$, avec $L'_{\cdot}$ plat \add{en tout degré}, qui est un quasi-isomorphisme. De ceci résulte formellement que $D^{-}(X)$ \uncertain{équivaut} à la catégorie obtenue en prenant (\uncertain{dans} la catégorie $K^{-}(X)^{\text{plat}}$ des complexes $L_{\cdot}$ bornés supérieurement à \uncertain{degrés} plats) \struck{\ill{}} en prenant un quotient par la \add{famille} \struck{relation} des quasi-isomorphismes.

Ceci va nous permettre de prendre des « foncteurs dérivés gauches » du foncteur $\otimes$, en prenant \uncertain{dans} $C^{-}(X)^{\text{plat}}$ le $\otimes$ ordinaire, \uncertain{qui} \uncertain{passe} au quotient par quasi-isomorphismes. Donc si $L_{\cdot}$, $M_{\cdot}$ sont deux complexes \uncertain{ordinaires} $\in\operatorname{Ob}D^{-}(X)$, on \struck{définit} calcule $L_{\cdot}\overset{\mathbb{L}}{\otimes}M_{\cdot}$ en prenant des résolutions plates $L'_{\cdot}\to L_{\cdot}$, $M'_{\cdot}\to M_{\cdot}$, et prenant $L'_{\cdot}\otimes M'_{\cdot}$.
\note{l'exposé du § 26 se poursuit au-delà de ce lot.}

\end{document}
